\documentclass{amsart}
% For dealing with references we use the comment environment
\usepackage{verbatim}
\newenvironment{reference}{\comment}{\endcomment}
%\newenvironment{reference}{}{}
\newenvironment{slogan}{\comment}{\endcomment}
\newenvironment{history}{\comment}{\endcomment}

% For commutative diagrams we use Xy-pic
\usepackage[all]{xy}

% We use 2cell for 2-commutative diagrams.
\xyoption{2cell}
\UseAllTwocells

% We use multicol for the list of chapters between chapters
\usepackage{multicol}

% This is generally recommended for better output
\usepackage{lmodern}
\usepackage[T1]{fontenc}

% For cross-file-references

% Package for hypertext links:
\usepackage{hyperref}

% For any local file, say "hello.tex" you want to link to please

% Theorem environments.
%
\theoremstyle{plain}
\newtheorem{theorem}[subsection]{Theorem}
\newtheorem{proposition}[subsection]{Proposition}
\newtheorem{lemma}[subsection]{Lemma}

\theoremstyle{definition}
\newtheorem{definition}[subsection]{Definition}
\newtheorem{example}[subsection]{Example}
\newtheorem{exercise}[subsection]{Exercise}
\newtheorem{situation}[subsection]{Situation}

\theoremstyle{remark}
\newtheorem{remark}[subsection]{Remark}
\newtheorem{remarks}[subsection]{Remarks}

\numberwithin{equation}{subsection}

% Macros
%
\def\lim{\mathop{\mathrm{lim}}\nolimits}
\def\colim{\mathop{\mathrm{colim}}\nolimits}
\def\Spec{\mathop{\mathrm{Spec}}}
\def\Hom{\mathop{\mathrm{Hom}}\nolimits}
\def\Ext{\mathop{\mathrm{Ext}}\nolimits}
\def\SheafHom{\mathop{\mathcal{H}\!\mathit{om}}\nolimits}
\def\SheafExt{\mathop{\mathcal{E}\!\mathit{xt}}\nolimits}
\def\Sch{\mathit{Sch}}
\def\Mor{\mathop{\mathrm{Mor}}\nolimits}
\def\Ob{\mathop{\mathrm{Ob}}\nolimits}
\def\Sh{\mathop{\mathit{Sh}}\nolimits}
\def\NL{\mathop{N\!L}\nolimits}
\def\CH{\mathop{\mathrm{CH}}\nolimits}
\def\proetale{{pro\text{-}\acute{e}tale}}
\def\etale{{\acute{e}tale}}
\def\QCoh{\mathit{QCoh}}
\def\Ker{\mathop{\mathrm{Ker}}}
\def\Im{\mathop{\mathrm{Im}}}
\def\Coker{\mathop{\mathrm{Coker}}}
\def\Coim{\mathop{\mathrm{Coim}}}

% Boxtimes
%
\DeclareMathSymbol{\boxtimes}{\mathbin}{AMSa}{"02}

%
% Macros for moduli stacks/spaces
%
\def\QCohstack{\mathcal{QC}\!\mathit{oh}}
\def\Cohstack{\mathcal{C}\!\mathit{oh}}
\def\Spacesstack{\mathcal{S}\!\mathit{paces}}
\def\Quotfunctor{\mathrm{Quot}}
\def\Hilbfunctor{\mathrm{Hilb}}
\def\Curvesstack{\mathcal{C}\!\mathit{urves}}
\def\Polarizedstack{\mathcal{P}\!\mathit{olarized}}
\def\Complexesstack{\mathcal{C}\!\mathit{omplexes}}
% \Pic is the operator that assigns to X its picard group, usage \Pic(X)
% \Picardstack_{X/B} denotes the Picard stack of X over B
% \Picardfunctor_{X/B} denotes the Picard functor of X over B
\def\Pic{\mathop{\mathrm{Pic}}\nolimits}
\def\Picardstack{\mathcal{P}\!\mathit{ic}}
\def\Picardfunctor{\mathrm{Pic}}
\def\Deformationcategory{\mathcal{D}\!\mathit{ef}}

\setlength{\emergencystretch}{3em}
\begin{document}
\title[Stacks Project, Tag 00RB]{114 --- Openness of fibrewise exactness for finite free complexes over flat finite-type algebras over a Noetherian base with Cohen-Macaulay fibres of fixed dimension}
\maketitle
\noindent Copyright (C) 2005--2025 Johan de Jong. Stacks Project authors. Source: \href{https://stacks.math.columbia.edu/tag/00RB}{Tag 00RB}.

\noindent Editorial difficulty note (not author text). Difficulty H2: The author spreads exactness from one localized fibre through determinantal rank bounds and converts the remaining grade conditions into regular-sequence conditions in the fibres. Cohen-Macaulay dimensional control makes these conditions open, giving a substantive family-wise exactness criterion beyond ordinary module exactness tests..

\noindent Editorial provenance note (not author text). History: excerpt from revision \texttt{a0444\allowbreak 6e57e\allowbreak c1fbc\allowbreak 252a8\allowbreak 71afc\allowbreak ec775\allowbreak 2fb28\allowbreak 07b14}, algebra.tex, lines 33306--33403. Reference presentation was adapted; original-author context and core support blocks were retained or added. The mathematical wording is unchanged. Licensed under GNU FDL 1.2 or later; no invariant sections or cover texts. See ../licenses/COPYING. Context blocks reproduce author definitions and supporting results; they are not additional counted problems.

\begin{situation}
\label{situation-complex}
Here $R$ is a ring, and we have a complex
$$
0
\to
R^{n_e}
\xrightarrow{\varphi_e}
R^{n_{e-1}}
\xrightarrow{\varphi_{e-1}}
\ldots
\xrightarrow{\varphi_{i + 1}}
R^{n_i}
\xrightarrow{\varphi_i}
R^{n_{i-1}}
\xrightarrow{\varphi_{i-1}}
\ldots
\xrightarrow{\varphi_1}
R^{n_0}
$$
In other words we require $\varphi_i \circ \varphi_{i + 1} = 0$
for $i = 1, \ldots, e - 1$.
\end{situation}

\begin{definition}
\label{definition-rank}
Let $R$ be a ring. Suppose that $\varphi : R^m \to R^n$ is a map
of finite free modules.
\begin{enumerate}
\item The {\it rank} of $\varphi$ is the maximal $r$ such that
$\wedge^r \varphi : \wedge^r R^m \to \wedge^r R^n$ is nonzero.
\item We let $I(\varphi) \subset R$ be the ideal generated by
the $r \times r$ minors of the matrix of $\varphi$, where $r$
is the rank as defined above.
\end{enumerate}
\end{definition}

\begin{definition}
\label{definition-relative-dimension}
Suppose that $R \to S$ is of finite type, and let
$\mathfrak q \subset S$ be a prime lying over a prime
$\mathfrak p$ of $R$.
We define the {\it relative dimension
of $S/R$ at $\mathfrak q$}, denoted
$\dim_{\mathfrak q}(S/R)$, to be the dimension
of $\Spec(S \otimes_R \kappa(\mathfrak p))$
at the point corresponding to $\mathfrak q$. We let
$\dim(S/R)$ be the supremum of $\dim_{\mathfrak q}(S/R)$
over all $\mathfrak q$. This is called the
{\it relative dimension of} $S/R$.
\end{definition}

\begin{proposition}
\label{proposition-what-exact}
\begin{reference}
\href{https://stacks.math.columbia.edu/bibliography}{Bibliography: WhatExact (Corollary 1)}
\end{reference}
In Situation \ref{situation-complex}, suppose $R$ is
a local Noetherian ring. The following are equivalent
\begin{enumerate}
\item $0 \to R^{n_e} \to R^{n_{e-1}} \to \ldots \to R^{n_0}$
is exact at $R^{n_e}, \ldots, R^{n_1}$, and
\item for all $i$, $1 \leq i \leq e$
the following two conditions are satisfied:
\begin{enumerate}
\item $\text{rank}(\varphi_i) = r_i$ where
$r_i = n_i - n_{i + 1} + \ldots + (-1)^{e-i-1} n_{e-1} + (-1)^{e-i} n_e$,
\item $I(\varphi_i) = R$, or $I(\varphi_i)$ contains a
regular sequence of length $i$.
\end{enumerate}
\end{enumerate}
\end{proposition}

\begin{proof}
If for some $i$ some matrix coefficient of $\varphi_i$
is not in $\mathfrak m$, then we apply Lemma \href{https://stacks.math.columbia.edu/tag/00MT}{lemma add trivial complex (Tag 00MT)}.
It is easy to see that the proposition for a complex and
for the same complex with a trivial complex added to it
are equivalent. Thus we may assume that all matrix entries
of each $\varphi_i$ are elements of the maximal ideal.
We may also assume that $e \geq 1$.

\medskip\noindent
Assume the complex is exact at $R^{n_e}, \ldots, R^{n_1}$.
Let $\mathfrak q \in \text{Ass}(R)$.
Note that the ring $R_{\mathfrak q}$ has depth $0$
and that the complex remains exact after localization at $\mathfrak q$.
We apply Lemmas \href{https://stacks.math.columbia.edu/tag/00MY}{lemma exact depth zero local (Tag 00MY)} and
\href{https://stacks.math.columbia.edu/tag/00MW}{lemma trivial case exact (Tag 00MW)} to the localized complex
over $R_{\mathfrak q}$. We conclude that
$\varphi_{i, \mathfrak q}$ has rank $r_i$ for all $i$.
Since $R \to \bigoplus_{\mathfrak q \in \text{Ass}(R)} R_\mathfrak q$
is injective (Lemma \href{https://stacks.math.columbia.edu/tag/0311}{lemma zero at ass zero (Tag 0311)}), we conclude that
$\varphi_i$ has rank $r_i$ over $R$ by the definition of rank as given
in Definition \ref{definition-rank}. Therefore we see that
$I(\varphi_i)_\mathfrak q = I(\varphi_{i, \mathfrak q})$
as the ranks do not change. Since all of the ideals
$I(\varphi_i)_{\mathfrak q}$, $e \geq i \geq 1$
are equal to $R_{\mathfrak q}$ (by the lemmas referenced above)
we conclude none of the ideals $I(\varphi_i)$ is contained in $\mathfrak q$.
This implies that $I(\varphi_e)I(\varphi_{e-1})\ldots I(\varphi_1)$
is not contained in any of the associated primes
of $R$. By Lemma \href{https://stacks.math.columbia.edu/tag/00DS}{lemma silly (Tag 00DS)} we may choose
$x \in I(\varphi_e)I(\varphi_{e - 1})\ldots I(\varphi_1)$,
$x \not \in \mathfrak q$ for all $\mathfrak q \in \text{Ass}(R)$.
Observe that $x$ is a nonzerodivisor (Lemma \href{https://stacks.math.columbia.edu/tag/00LD}{lemma ass zero divisors (Tag 00LD)}).
According to Lemma \href{https://stacks.math.columbia.edu/tag/00MZ}{lemma div x exact one less (Tag 00MZ)}
the complex $0 \to (R/xR)^{n_e} \to \ldots \to (R/xR)^{n_1}$ is exact
at $(R/xR)^{n_e}, \ldots, (R/xR)^{n_2}$. By induction
on $e$ all the ideals $I(\varphi_i)/xR$ have a regular
sequence of length $i - 1$. This proves that $I(\varphi_i)$
contains a regular sequence of length $i$.

\medskip\noindent
Assume (2)(a) and (2)(b) hold. We will prove that (1) holds by
induction on $\dim(R)$. If $\dim(R) = 0$, then we must have
$I(\varphi_i) = R$ for $1 \leq i \leq e$ by (2)(b). Since the
coefficients of $\varphi_i$ are contained in the maximal ideal
this can happen only if $r_i = 0$ for all $i$. By (2)(a) we
conclude that $e = 0$ and (1) holds. Assume $\dim(R) > 0$.
We claim that for any prime $\mathfrak p \subset R$ conditions
(2)(a) and (2)(b) hold for the complex
$0 \to R_\mathfrak p^{n_e} \to R_\mathfrak p^{n_{e - 1}} \to \ldots \to
R_\mathfrak p^{n_0}$ with maps $\varphi_{i, \mathfrak p}$
over $R_\mathfrak p$. Namely, since $I(\varphi_i)$ contains a
nonzero divisor, the image of $I(\varphi_i)$ in $R_\mathfrak p$
is nonzero. This implies that the rank of $\varphi_{i, \mathfrak p}$
is the same as the rank of $\varphi_i$: the rank as defined above
of a matrix $\varphi$ over a ring $R$ can only drop when passing
to an $R$-algebra $R'$ and this happens if and only if $I(\varphi)$
maps to zero in $R'$. Thus (2)(a) holds. Having said this
we know that $I(\varphi_{i, \mathfrak p}) = I(\varphi_i)_\mathfrak p$
and we see that (2)(b) is preserved under localization as well.
By induction on the dimension of $R$ we may assume the complex
is exact when localized at any nonmaximal prime $\mathfrak p$ of $R$.
Thus $\Ker(\varphi_i)/\Im(\varphi_{i + 1})$ has support contained in
$\{\mathfrak m\}$ and hence if nonzero has depth $0$.
As $I(\varphi_i) \subset \mathfrak m$ for all $i$ because
of what was said in the first paragraph of the proof, we
see that (2)(b) implies $\text{depth}(R) \geq e$.
By Lemma \href{https://stacks.math.columbia.edu/tag/00N0}{lemma acyclic (Tag 00N0)} we see
that the complex is exact at $R^{n_e}, \ldots, R^{n_1}$
concluding the proof.
\end{proof}


\begin{lemma}
\label{lemma-open-regular-sequence}
Let $R \to S$ be a finite type ring map. Let $d$ be an integer such that all
fibres $S \otimes_R \kappa(\mathfrak p)$ are Cohen-Macaulay and equidimensional
of dimension $d$. Let $f_1, \ldots, f_i$ be elements of $S$. The set
$$
\{ \mathfrak q \in V(f_1, \ldots, f_i)
\mid f_1, \ldots, f_i
\text{ are a regular sequence in }
S_{\mathfrak q}/\mathfrak p S_{\mathfrak q}
\text{ where }\mathfrak p = R \cap \mathfrak q
\}
$$
is open in $V(f_1, \ldots, f_i)$.
\end{lemma}

\begin{proof}
Write $\overline{S} = S/(f_1, \ldots, f_i)$.
Suppose $\mathfrak q$ is an element of the set defined in the
lemma, and $\mathfrak p$ is the corresponding prime of $R$.
We will use relative dimension as defined in
Definition \ref{definition-relative-dimension}.
First, note that $d = \dim_{\mathfrak q}(S/R) =
\dim(S_{\mathfrak q}/\mathfrak pS_{\mathfrak q}) +
\text{trdeg}_{\kappa(\mathfrak p)}\ \kappa(\mathfrak q)$
by Lemma \href{https://stacks.math.columbia.edu/tag/00P1}{lemma dimension at a point finite type field (Tag 00P1)}.
Since $f_1, \ldots, f_i$ form a regular sequence in the
Noetherian local ring $S_{\mathfrak q}/\mathfrak pS_{\mathfrak q}$
Lemma \href{https://stacks.math.columbia.edu/tag/00KW}{lemma one equation (Tag 00KW)} tells us that
$\dim(\overline{S}_{\mathfrak q}/\mathfrak p\overline{S}_{\mathfrak q})
= \dim(S_{\mathfrak q}/\mathfrak pS_{\mathfrak q}) - i$.
We conclude that $\dim_{\mathfrak q}(\overline{S}/R)
= \dim(\overline{S}_{\mathfrak q}/\mathfrak p\overline{S}_{\mathfrak q}) +
\text{trdeg}_{\kappa(\mathfrak p)}\ \kappa(\mathfrak q)
= d - i$ by Lemma \href{https://stacks.math.columbia.edu/tag/00P1}{lemma dimension at a point finite type field (Tag 00P1)}.
By Lemma \href{https://stacks.math.columbia.edu/tag/00QH}{lemma dimension fibres bounded open upstairs (Tag 00QH)}
we have $\dim_{\mathfrak q'}(\overline{S}/R) \leq d - i$
for all $\mathfrak q' \in V(f_1, \ldots, f_i) = \Spec(\overline{S})$
in a neighbourhood of $\mathfrak q$. Thus after replacing
$S$ by $S_g$ for some $g \in S$, $g \not \in \mathfrak q$
we may assume that the inequality holds for all
$\mathfrak q'$. The result follows from Lemma
\href{https://stacks.math.columbia.edu/tag/00R9}{lemma CM dim finite type (Tag 00R9)}.
\end{proof}


\begin{lemma}
\label{lemma-exact-on-fibres-open}
Let $R \to S$ be a ring map. Consider a finite homological complex of
finite free $S$-modules:
$$
F_{\bullet} :
0
\to
S^{n_e}
\xrightarrow{\varphi_e}
S^{n_{e-1}}
\xrightarrow{\varphi_{e-1}}
\ldots
\xrightarrow{\varphi_{i + 1}}
S^{n_i}
\xrightarrow{\varphi_i}
S^{n_{i-1}}
\xrightarrow{\varphi_{i-1}}
\ldots
\xrightarrow{\varphi_1}
S^{n_0}
$$
For every prime $\mathfrak q$ of $S$ consider the
complex $\overline{F}_{\bullet, \mathfrak q} =
F_{\bullet, \mathfrak q} \otimes_R \kappa(\mathfrak p)$
where $\mathfrak p$ is inverse image of $\mathfrak q$ in $R$.
Assume $R$ is Noetherian and there exists an integer $d$ such
that $R \to S$ is finite type, flat
with fibres $S \otimes_R \kappa(\mathfrak p)$
Cohen-Macaulay of dimension $d$.
The set
$$
\{\mathfrak q \in \Spec(S) \mid
\overline{F}_{\bullet, \mathfrak q}\text{ is exact}\}
$$
is open in $\Spec(S)$.
\end{lemma}

\begin{proof}
Let $\mathfrak q$ be an element of the set defined in the lemma.
We are going to use Proposition \ref{proposition-what-exact}
to show there exists a $g \in S$, $g \not \in \mathfrak q$
such that $D(g)$ is contained in the set defined in the lemma.
In other words, we are going to show that after replacing $S$
by $S_g$, the set of the lemma is all of $\Spec(S)$.
Thus during the proof we will, finitely often, replace
$S$ by such a localization.
Recall that Proposition \ref{proposition-what-exact}
characterizes exactness of complexes
in terms of ranks of the maps $\varphi_i$ and the ideals
$I(\varphi_i)$, in case the ring is local. We first address
the rank condition. Set
$r_i = n_i - n_{i + 1} + \ldots + (-1)^{e - i} n_e$.
Note that $r_i + r_{i + 1} = n_i$ and note that
$r_i$ is the expected rank of $\varphi_i$ (in the
exact case).

\medskip\noindent
By Lemma \href{https://stacks.math.columbia.edu/tag/00MI}{lemma complex exact mod (Tag 00MI)} we see that if
$\overline{F}_{\bullet, \mathfrak q}$ is exact, then
the localization $F_{\bullet, \mathfrak q}$ is exact.
In particular the complex $F_\bullet$ becomes
exact after localizing by an element
$g \in S$, $g \not \in \mathfrak q$. In this case
Proposition \ref{proposition-what-exact} applied
to all localizations of $S$ at prime ideals
implies that all $(r_i + 1) \times (r_i + 1)$-minors
of $\varphi_i$ are zero. Thus we see that the rank
of $\varphi_i$ is at most $r_i$.

\medskip\noindent
Let $I_i \subset S$ denote the ideal generated
by the $r_i \times r_i$-minors of the matrix
of $\varphi_i$. By Proposition \ref{proposition-what-exact}
the complex $\overline{F}_{\bullet, \mathfrak q}$ is exact
if and only if for every $1 \leq i \leq e$ we have
either $(I_i)_{\mathfrak q} = S_{\mathfrak q}$ or
$(I_i)_{\mathfrak q}$ contains a $S_{\mathfrak q}/\mathfrak p
S_{\mathfrak q}$-regular sequence of length $i$.
Namely, by our choice of $r_i$ above and by the
bound on the ranks of the $\varphi_i$ this is the
only way the conditions of Proposition \ref{proposition-what-exact}
can be satisfied.

\medskip\noindent
If $(I_i)_{\mathfrak q} = S_{\mathfrak q}$, then after localizing $S$ at
some element $g \not\in \mathfrak q$ we may assume that
$I_i = S$. Clearly, this is an open condition.

\medskip\noindent
If $(I_i)_{\mathfrak q} \not = S_{\mathfrak q}$, then we have
a sequence $f_1, \ldots, f_i \in (I_i)_{\mathfrak q}$ which
form a regular sequence in $S_{\mathfrak q}/\mathfrak pS_{\mathfrak q}$.
Note that for any prime $\mathfrak q' \subset S$ such that
$(f_1, \ldots, f_i) \not \subset \mathfrak q'$ we have
$(I_i)_{\mathfrak q'} = S_{\mathfrak q'}$.
Thus the result follows from Lemma \ref{lemma-open-regular-sequence}.
\end{proof}
\end{document}
